\chapter{Introdução}
\label{cap:introducao}

O Sistema Integrado de Gestão de Atividades Acadêmicas (SIGAA), desenvolvido pela Universidade Federal do Rio Grande do Norte (UFRN) e adotado por diversas universidades federais, é o sistema pelo qual o estudante de graduação da Universidade de Brasília (UnB) consulta notas, horários, histórico e índices acadêmicos, emite documentos, realiza matrícula e acessa as turmas virtuais \cite{sigaa}. Não há alternativa institucional: uma dificuldade no sistema se converte diretamente em atraso, erro de planejamento ou dependência de colegas e de ferramentas de terceiros.

Há indícios de que o SIGAA não atende bem esse público. A interface foi concebida para computador de mesa, com largura fixa e menus que se abrem pela passagem do ponteiro do mouse; informações essenciais, como o horário das turmas, aparecem em notação codificada (por exemplo, ``24M34''); e estudantes recorrem a ferramentas criadas pela própria comunidade para tarefas que o sistema não resolve de forma satisfatória, como o planejamento do fluxo de disciplinas até a formatura. Esses indícios, porém, são percepções. Para que possam apoiar decisões de melhoria, precisam ser transformados em evidências mensuráveis e comparadas a critérios definidos antes da coleta.

A família de normas SQuaRE (\textit{Systems and software Quality Requirements and Evaluation}) oferece o arcabouço para essa transformação: um modelo de qualidade de produto \cite{iso25010_2023}, um catálogo de medidas \cite{iso25023} e um processo de avaliação em fases \cite{iso25040}. O método \textit{Goal-Question-Metric} (GQM) liga os objetivos da avaliação às questões e às métricas que os respondem \cite{basili1994}, e o \textit{Practical Software and Systems Measurement} (PSM) orienta a medição do próprio projeto para apoiar decisões de gestão \cite{mcgarry2002}.

\section{Objetivos}
\label{sec:objetivos}

O objetivo geral deste trabalho é \textbf{avaliar a qualidade do portal do discente de graduação do SIGAA/UnB, do ponto de vista do estudante}, para identificar com evidências os pontos fracos do sistema nas tarefas acadêmicas mais frequentes e apoiar a priorização de melhorias pela Secretaria de Tecnologia da Informação (STI/UnB). Os objetivos específicos são:

\begin{enumerate}
	\item estabelecer os requisitos da avaliação: propósito, partes interessadas, produto, modelo de qualidade, características priorizadas e escopo (Fase 1);
	\item especificar a avaliação pelo método GQM, com métricas, níveis de pontuação e critérios de julgamento definidos antes da coleta (Fase 2);
	\item projetar e executar a avaliação, julgar os resultados e discutir as relações entre as características (Fases 3 e 4);
	\item executar ao menos uma ação de melhoria justificada pelas evidências obtidas (Fase 4).
\end{enumerate}

Este artigo corresponde à Entrega 1 do projeto e cobre os dois primeiros objetivos, conforme o roteiro da disciplina \cite{ramos2026roteiro}.

\section{Organização do texto e repositório}
\label{sec:organizacao}

O \autoref{cap:referencial} apresenta o referencial teórico; o \autoref{cap:metodologia}, a metodologia de desenvolvimento do projeto, incluindo a gestão com indicadores PSM/CID; o \autoref{cap:fase1} descreve a Fase 1 e o \autoref{cap:fase2}, a Fase 2; o \autoref{cap:consideracoes} traz as considerações parciais e os próximos passos. A tabela de contribuição da equipe está no \autoref{ap:contribuicao}.

Todos os artefatos citados (páginas detalhadas de cada fase, dados, instrumentos de coleta e histórico de versões) estão no repositório da equipe, disponível em \urlrepositorio, e publicados na página do projeto (GitPage), em \urlgitpage.
